%%**************************************************************
%% Vorlage fuer Bachelorarbeiten (o.ä.) der DHBW
%%
%% Autor: Tobias Dreher, Yves Fischer
%% Datum: 06.07.2011
%%
%% Autor: Michael Gruben
%% Datum: 15.05.2013
%%
%% Geaendert 2026: lokaler Build mit latexmk, VS-Code-Live-Vorschau,
%% Sprachschalter, alle Angaben in einem Block, Bugfixes.
%% Lizenz: Apache 2.0 (siehe LICENSE). Anleitung: CLAUDE.md im Hauptordner.
%%**************************************************************
\documentclass[%
	enabledeprecatedfontcommands, % agsm.bst schreibt \bf ins Literaturverzeichnis
	pdftex,
	oneside,		% Einseitiger Druck.
	12pt,			% Schriftgroesse
	parskip=half,	% Halbe Zeile Abstand zwischen Absätzen.
	headsepline,	% Linie nach Kopfzeile.
	footsepline,	% Linie vor Fusszeile.
	abstract=true,	% Überschrift über dem Abstract
	ngerman,		% Translator
]{scrreprt}

%% ==============================================================
%%  DEINE ANGABEN / YOUR DATA -- nur dieser Block / only this block
%%  Leeres Feld = Zeile entfällt auf dem Deckblatt / empty field =
%%  row dropped from the title page, e.g. \newcommand{\gutachter}{}
%%  & immer als \& schreiben / always write & as \&
%% ==============================================================
% Sprache / language: einmal mit ./set-language.sh de|en wählen /
% choose once with ./set-language.sh de|en (sets this switch and the skeleton)
\newif\ifenglisch
\englischfalse
\newif\ifsperrvermerk
\sperrvermerktrue   % \sperrvermerkfalse = ohne Sperrvermerk / no confidentiality page

\newcommand{\titel}{Titel der Arbeit}                 % Titel / title
\newcommand{\untertitel}{Untertitel der Arbeit}       % Untertitel / subtitle (optional)
\newcommand{\arbeit}{Projektarbeit (T2000)}           % Art der Arbeit / type of work
\newcommand{\autor}{Vorname Nachname}                 % Name / name
\newcommand{\matrikelnr}{1234567}                     % Matrikelnummer / student ID
\newcommand{\kurs}{TINF24A}                           % Kurs / course
\newcommand{\studiengang}{Informatik}                 % Studiengang / study program
\newcommand{\studienrichtung}{Angewandte Informatik}  % Studienrichtung / specialization (optional)
\newcommand{\standort}{Mannheim}                      % DHBW-Standort / DHBW campus
\newcommand{\firma}{Musterfirma GmbH \& Co. KG}       % Ausbildungsfirma / partner company
\newcommand{\firmenort}{Musterstraße 1, 12345 Musterstadt} % Adresse / address
\newcommand{\betreuer}{Vorname Nachname}              % Betreuer Firma / in-company supervisor
\newcommand{\gutachter}{Prof. Dr. Vorname Nachname}   % Gutachter DHBW / university reviewer
\newcommand{\zeitraum}{12 Wochen}                     % Bearbeitungszeitraum / processing period
\newcommand{\datumAbgabe}{01.09.2026}                 % Abgabedatum / submission date
\newcommand{\abgabeort}{Mannheim}                     % Abgabeort / place of submission
%% ==============================================================

%!TEX root = ../dokumentation.tex
%
% Alle Stil-Einstellungen der Arbeit, in Abschnitten:
%    1 Seite und Zeilen          6 Quellcode
%    2 Schrift                   7 Sonstiges (Fußnoten, Entwurfsnotizen)
%    3 Sprache                   8 Links und PDF-Infos (hyperref)
%    4 Farben                    9 Zitierweise
%    5 Abbildungen, Tabellen,   10 Glossar und Abkürzungen
%      Diagramme (TikZ-Stile)
% Persönliche Angaben stehen nicht hier, sondern oben in dokumentation.tex.
% Reihenfolge beachten: hyperref (8) lädt spät; 9 und 10 brauchen es.
%

%% ---- 1 Seite und Zeilen --------------------------------------------
% Seitenränder: fullpage = rundum ca. 2,5 cm. Feste Maße stattdessen:
%   \usepackage[left=3cm,right=2.5cm,top=2.5cm,bottom=2.5cm]{geometry}
\usepackage{fullpage}

% Zeilenabstand: onehalfspacing = 1,5-zeilig; singlespacing, doublespacing
% oder freier Wert mit \setstretch{1.3} nach dem Laden.
\usepackage[onehalfspacing]{setspace}

% Hurenkinder und Schusterjungen verhindern
% http://projekte.dante.de/DanteFAQ/Silbentrennung
\clubpenalty=10000
\widowpenalty=10000
\displaywidowpenalty=10000

%% ---- 2 Schrift -----------------------------------------------------
\usepackage[utf8]{inputenc}
\usepackage[T1]{fontenc}

% Schriftart Palatino. Alternativen (eine Zeile statt palatino):
%   \usepackage{lmodern}               Latin Modern (LaTeX-Standard)
%   \usepackage{libertine}             Linux Libertine
%   \usepackage{mathptmx}              Times
%   \usepackage{newpxtext,newpxmath}   Palatino-Klon mit passender Mathe
\usepackage{palatino}

% Zeilenumbruch und mehr
\usepackage[activate]{microtype}

% Überschriften setzt KOMA-Script serifenlos fett. In der Textschrift:
%\addtokomafont{disposition}{\rmfamily}

%% ---- 3 Sprache -----------------------------------------------------
% Die Arbeit hat genau eine Sprache: ./set-language.sh de|en setzt
% \englischtrue / \englischfalse in dokumentation.tex und das passende Gerüst.
% Deckblatt, Erklärung, Verzeichnisse, Zitate und Verweise folgen dem Schalter.
% \deen{deutsch}{english}: Text je nach Sprache, für Vorlagenteile wie das Deckblatt
\newcommand{\deen}[2]{\ifenglisch#2\else#1\fi}
% main= explizit: 'ngerman' steht auch in den Klassenoptionen und wird
% global durchgereicht, wodurch die reine Reihenfolge nicht mehr entscheidet.
\ifenglisch
	\usepackage[ngerman,main=english]{babel}
\else
	\usepackage[english,main=ngerman]{babel}
\fi

% Anführungszeichen: \enquote{...} setzt sie passend zur Sprache
\usepackage[babel,german=quotes]{csquotes}

% Überschriften, die babel nicht kennt
\ifenglisch
	\newcommand{\abkuerzungstitel}{List of Abbreviations}
	\newcommand{\glossartitel}{Glossary}
\else
	\newcommand{\abkuerzungstitel}{Abkürzungsverzeichnis}
	\newcommand{\glossartitel}{Glossar}
\fi

% Index-Erstellung
\usepackage{makeidx}

%% ---- 4 Farben ------------------------------------------------------
% xcolor, not color: listings and TikZ load xcolor later, and definitions made
% with plain `color` then raise "Incompatible color definition".
\usepackage{xcolor}
% Farbe aller Links im PDF. Für den Druck schwarz: {rgb}{0,0,0}
\definecolor{LinkColor}{rgb}{0,0,0.2}
% Grundfarbe aller TikZ-Diagramme (Abschnitt 5), z.B. blue!60!black
\colorlet{akzent}{black}

%% ---- 5 Abbildungen, Tabellen, Diagramme ----------------------------
% Grafiken liegen in images/; Dateiname ohne Endung angeben.
\usepackage{graphicx}
\graphicspath{{images/}}

% Gleitumgebungen (Bilder, Tabellen, usw\ldots) lassen sich mit H an genau der
% definierten Stelle platzieren
\usepackage{float}

% für textumflossene Grafiken
\usepackage{wrapfig}

% Tabellen: booktabs (\toprule, \midrule, \bottomrule), tabularx (Spalten
% füllen die Breite), array (vertikale Platzierung), longtable (über Seiten).
\usepackage{booktabs}
\usepackage{tabularx}
\usepackage{array}
\usepackage{longtable}
% Gilt für ALLE Tabellen, auch das Deckblatt: Spaltenabstand und Zeilenhöhe
\setlength{\tabcolsep}{10pt}
\renewcommand{\arraystretch}{1.5}

% Formeln
\usepackage{amsmath}
%\usepackage{amssymb}

% Pakete um Textteile drehen zu können, oder eine Seite Querformat anzeigen kann.
%\usepackage{rotating}
%\usepackage{lscape}

% TikZ-Diagramme mit einem gemeinsamen Stilbaukasten, damit alle Diagramme
% gleich aussehen. Farbe zentral über 'akzent' (Abschnitt 4).
%   flow   = Grundstil fürs Bild:  \begin{tikzpicture}[flow]
%   fio    = Ein-/Ausgabe (gerundet, leicht gefüllt)
%   fstage = Verarbeitungsschritt
%   fkey   = hervorgehobener Schritt (dunkler gefüllt)
%   fmuted = Nebenschritt (heller Rahmen, graue Schrift)
%   fgate  = Entscheidung (Raute)
%   fedge  = Pfeil;  fcond = bedingter Pfeil (gestrichelt)
%   flab   = Beschriftung an einem Pfeil;  fnote = kleine Randnotiz
% Beispiel: content/beispiele.tex
\usepackage{tikz}
\usetikzlibrary{arrows.meta,positioning,shapes.geometric,calc}
% Diagramme nie breiter als der Satzspiegel: max width verkleinert nur,
% vergroessert aber nie (im Gegensatz zu \resizebox).
\usepackage{adjustbox}
\tikzset{
  flow/.style={font=\footnotesize\sffamily, line cap=round},
  fio/.style={draw=akzent!70, fill=akzent!8, rounded corners=2.2mm, inner xsep=2.4mm, inner ysep=1.5mm, align=center},
  fstage/.style={draw=akzent, line width=0.7pt, fill=white, rounded corners=0.6mm, inner xsep=2.4mm, inner ysep=1.6mm, align=center},
  fkey/.style={fstage, fill=akzent!15},
  fmuted/.style={draw=akzent!35, fill=white, rounded corners=0.6mm, inner xsep=2.2mm, inner ysep=1.5mm, align=center, text=black!60},
  fgate/.style={draw=akzent!50, fill=akzent!4, diamond, aspect=2.4, inner sep=0.8mm, align=center},
  fedge/.style={-{Stealth[length=2mm]}, draw=akzent!55, semithick},
  fcond/.style={fedge, dashed},
  flab/.style={font=\scriptsize\sffamily, text=black!60, inner sep=1pt, fill=white},
  fnote/.style={font=\scriptsize\sffamily, text=black!45, inner sep=1pt},
}

%% ---- 6 Quellcode ---------------------------------------------------
% listings statt minted: braucht kein -shell-escape und kein Python.
% Sprache pro Listing ändern: \begin{lstlisting}[language=Java]
\usepackage{listings}
\definecolor{noteframe}{gray}{0.55}
\lstset{
  basicstyle=\ttfamily\scriptsize,
  keywordstyle=\color{black}\bfseries,
  commentstyle=\itshape\color{noteframe},
  stringstyle=\color{black},
  numbers=left, numberstyle=\tiny\color{noteframe}, numbersep=7pt,
  frame=leftline, framesep=6pt, rulecolor=\color{noteframe},
  breaklines=true, breakatwhitespace=false, breakindent=12pt,
  showstringspaces=false, tabsize=2, columns=fullflexible,
  language=Python, upquote=true, aboveskip=10pt, belowskip=10pt,
  captionpos=b, xleftmargin=10pt,
  % listings reads its body byte by byte, so every non-ASCII character that
  % appears in the source code needs an explicit mapping.
  extendedchars=true, inputencoding=utf8,
  literate={ä}{{\"a}}1 {ö}{{\"o}}1 {ü}{{\"u}}1
           {Ä}{{\"A}}1 {Ö}{{\"O}}1 {Ü}{{\"U}}1
           {ß}{{\ss}}1 {–}{{--}}1 {—}{{---}}1
           {─}{{-}}1,
}
\renewcommand{\lstlistingname}{Listing}
\ifenglisch\else\renewcommand{\lstlistlistingname}{Quellcodeverzeichnis}\fi

%% ---- 7 Sonstiges ---------------------------------------------------
% Fussnoten
\usepackage[perpage, hang, multiple, stable]{footmisc}

% für die Darstellung des Euro-Symbols: \euro{}
\usepackage[right]{eurosym}

% Entwurfsnotizen: \begin{notiz} ... \end{notiz} wird rot gedruckt.
% Vor der Abgabe \excludecomment{notiz} einkommentieren: alle Notizen verschwinden.
% (Nicht 'k' nennen: \k ist der Ogonek-Akzent für ą/ę, den hyperref neu definiert.)
\usepackage{comment}
% Zwischendatei als dokumentation.cut statt comment.cut, damit latexmk -c sie aufräumt
\def\DefaultCutFileName{\def\CommentCutFile{\jobname.cut}}
\specialcomment{notiz}{\begingroup\color{red}}{\endgroup}
%\excludecomment{notiz}

%% ---- 8 Links und PDF-Infos -----------------------------------------
\usepackage[%
	pdftitle={\titel},
	pdfauthor={\autor},
	pdfsubject={\arbeit},
	pdfcreator={pdflatex, LaTeX with KOMA-Script},
	pdfpagemode=UseOutlines, % Beim Oeffnen Inhaltsverzeichnis anzeigen
	pdfdisplaydoctitle=true, % Dokumenttitel statt Dateiname anzeigen.
]{hyperref}

% (Farb-)einstellungen für die Links im PDF.
% Ganz ohne farbige Links: colorlinks=false, hidelinks
\hypersetup{%
	colorlinks=true, % Aktivieren von farbigen Links im Dokument
	linkcolor=LinkColor, % Farbe festlegen
	citecolor=LinkColor,
	filecolor=LinkColor,
	menucolor=LinkColor,
	urlcolor=LinkColor,
	bookmarksnumbered=true % Überschriftsnummerierung im PDF Inhalt anzeigen.
}
\ifenglisch\hypersetup{pdflang=en}\else\hypersetup{pdflang=de}\fi

% \autoref{lst:...} schreibt "Listing" vor die Nummer
\providecommand*{\lstlistingautorefname}{Listing}
% Englisch: hyperref schreibt "chapter 2"; großgeschrieben passt es auch am Satzanfang
\ifenglisch
	\AtBeginDocument{%
		\def\chapterautorefname{Chapter}%
		\def\sectionautorefname{Section}%
		\def\subsectionautorefname{Section}%
		\def\subsubsectionautorefname{Section}%
		\def\appendixautorefname{Appendix}%
	}
\fi

%nur ein latex-Durchlauf für die Aktualisierung von Verzeichnissen nötig
\usepackage{bookmark}

%% ---- 9 Zitierweise -------------------------------------------------
% Literaturverweise nach Harvard (Autor, Jahr) mit agsm.bst
\usepackage[dcucite]{harvard}
\ifenglisch\renewcommand{\harvardand}{and}\else\renewcommand{\harvardand}{und}\fi
% agsm.bst emits \harvardurl{<raw url>}; the original grabs its argument
% normally, so any '_' in a URL raises "Missing $ inserted".
% Handing the braces straight to \url makes it read them verbatim.
\AtBeginDocument{\renewcommand{\harvardurl}{\textbf{URL:} \url}}
% Ab drei Autoren immer (Autor et al., Jahr)
\citationmode{abbr}

% Zitierbefehle (natbib-Namen; harvard.sty kennt dieselben Formen unter
% anderem Namen):
%   \citep{key}          -> (Autor, Jahr)        == \cite{key}
%   \citep[S.~12]{key}   -> (Autor, Jahr, S. 12)
%   \citet{key}          -> Autor (Jahr)         == \citeasnoun{key}
\newcommand{\citep}[2][]{\if\relax\detokenize{#1}\relax\cite{#2}\else\cite[#1]{#2}\fi}
\newcommand{\citet}[2][]{\if\relax\detokenize{#1}\relax\citeasnoun{#2}\else\citeasnoun[#1]{#2}\fi}

%% ---- 10 Glossar und Abkürzungen ------------------------------------
% Glossar
\usepackage[
	nonumberlist, %keine Seitenzahlen anzeigen
	%acronym,      %ein Abkürzungsverzeichnis erstellen
	%section,     %im Inhaltsverzeichnis auf section-Ebene erscheinen
	toc,          %Einträge im Inhaltsverzeichnis
]{glossaries}

% Akronyme: \ac{API} -- beim ersten Mal ausgeschrieben (als Fußnote)
\usepackage[printonlyused,footnote]{acronym}


\makeglossaries
%!TEX root = ../dokumentation.tex

%
% Glossareintraege --> referenz, name, beschreibung
% Aufruf im Text mit \gls{...}; \glsaddall in dokumentation.tex sorgt dafuer,
% dass alle Eintraege gedruckt werden, auch die im Text nicht verlinkten.
%

\newglossaryentry{synctex}{name={SyncTeX},
  description={Verknüpfung zwischen Quelltext und PDF: Ein Sprung im einen
  landet an der passenden Stelle im anderen}}

\newglossaryentry{latexmk}{name={latexmk},
  description={Build-Werkzeug, das pdflatex, BibTeX und makeglossaries so oft
  aufruft, bis alle Verweise aufgelöst sind}}


\begin{document}

	% Deckblatt
	\begin{spacing}{1}
		%!TEX root = ../dokumentation.tex
% Deckblatt. Alle Inhalte kommen aus dem Angaben-Block in dokumentation.tex.
% Ein leeres Feld lässt seine Zeile weg.

% \deckblock{Abstand}{Format oder Vortext}{\feld}: zentrierte Zeile
\newcommand{\deckblock}[3]{\if\relax\detokenize\expandafter{#3}\relax\else\vspace*{#1}{#2#3}\\\fi}
% \deckzeile{Bezeichnung}{\feld}: Zeile der Angaben-Tabelle unten
\newcommand{\deckzeile}[2]{\if\relax\detokenize\expandafter{#2}\relax\else\textbf{#1} \> #2\\\fi}

\begin{titlepage}
	% Logos: Firmenlogo links (images/logo.png, ohne Datei kein Logo), DHBW rechts
	\begin{longtable}{p{.4\textwidth} p{.4\textwidth}}
	  \IfFileExists{images/logo.png}{\includegraphics[height=2.6cm]{images/logo.png}}{} &
	  \includegraphics[height=2.6cm]{images/dhbw.png}
	\end{longtable}
	\enlargethispage{20mm}
	\begin{center}
		\vspace*{12mm}	Duale Hochschule Baden-Württemberg \standort\\
		\vspace*{12mm}	{\LARGE\bfseries \titel}\\
		\deckblock{4mm}{\large}{\untertitel}
		\vspace*{12mm}	{\large\bfseries \arbeit}\\
		\deckblock{12mm}{Studiengang~}{\studiengang}
		\deckblock{3mm}{Studienrichtung~}{\studienrichtung}
		\vspace*{3mm}	an der Dualen Hochschule Baden-Württemberg \standort\\
		\vspace*{12mm}	von\\
		\vspace*{3mm}	{\large\bfseries \autor}\\
		\vspace*{12mm}	\datumAbgabe\\
	\end{center}
	\vfill
	\begin{spacing}{1.2}
	\begin{tabbing}
		\textbf{Gutachter der Dualen Hochschule}\quad \= \kill
		\deckzeile{Bearbeitungszeitraum}{\zeitraum}
		\textbf{Matrikelnummer, Kurs} \> \matrikelnr, \kurs\\
		\deckzeile{Ausbildungsfirma}{\firma}
		\deckzeile{}{\firmenort}
		\deckzeile{Betreuer der Ausbildungsfirma}{\betreuer}
		\deckzeile{Gutachter der Dualen Hochschule}{\gutachter}
	\end{tabbing}
	\end{spacing}
\end{titlepage}

	\end{spacing}
	\newpage

	\renewcommand{\thepage}{\Roman{page}}
	\setcounter{page}{1}

	% Sperrvermerk (abschalten oben mit \sperrvermerkfalse)
	\ifsperrvermerk
		%!TEX root = ../dokumentation.tex
% Sperrvermerk direkt hinter der Titelseite.
% Abschalten: \sperrvermerkfalse in dokumentation.tex.

\thispagestyle{empty}

\ifenglisch
\section*{Confidentiality Notice}

\vspace*{2em}

The present work is based on internal, confidential data and information of \firma.
The content of this work may not, in whole or in part, be made accessible to persons outside the examination and evaluation process unless explicit permission has been granted by the dual partner.
\else
\section*{Sperrvermerk}

\vspace*{2em}

Die vorliegende Arbeit basiert auf internen, vertraulichen Daten und Informationen des Unternehmens \firma.
Der Inhalt dieser Arbeit darf weder als Ganzes noch in Auszügen Personen außerhalb des Prüfungs- und Evaluationsverfahrens zugänglich gemacht werden, sofern keine anders lautende Genehmigung des Dualen Partners vorliegt.
\fi

		\newpage
	\fi

	% Erklärung
	%!TEX root = ../dokumentation.tex
% Eidesstattliche Erklärung. Titel, Art der Arbeit, Ort und Datum kommen aus
% dokumentation.tex. Den Wortlaut vor der Abgabe mit der Prüfungsordnung bzw.
% den Vorgaben des Studiengangs abgleichen; manche verlangen zusätzlich eine
% Erklärung zur Nutzung von KI-Werkzeugen.

\thispagestyle{empty}

\ifenglisch
\section*{Declaration}

\vspace*{2em}

I hereby declare that I have written the present work with the title \enquote{\titel} independently and that I have used no sources or aids other than those indicated.
I further declare that the submitted electronic version is identical to the printed version.
\else
\section*{Erklärung}

\vspace*{2em}

Ich versichere hiermit, dass ich meine \arbeit{} mit dem Thema: \enquote{\titel} selbstständig verfasst und keine anderen als die angegebenen Quellen und Hilfsmittel benutzt habe.
Ich versichere zudem, dass die eingereichte elektronische Fassung mit der gedruckten Fassung übereinstimmt.
\fi

\vspace{3em}

\abgabeort, \datumAbgabe
\vspace{4em}

\autor

	\newpage

	% Abstract
	%!TEX root = ../dokumentation.tex
% Abstract: etwa eine halbe Seite, ein Satz pro Zeile.
% Zusätzliche deutsche Kurzfassung: diesen Block kopieren und davor
% \renewcommand{\abstractname}{Kurzfassung} setzen.

\pagestyle{empty}

\renewcommand{\abstractname}{Abstract}
\begin{abstract}
Hier steht die Zusammenfassung der Arbeit.
Sie nennt das Problem, das Vorgehen, die wichtigsten Ergebnisse und die Schlussfolgerung.
Auch hier gilt: ein Satz pro Zeile.
\end{abstract}

	\newpage

	\pagestyle{plain}

	% Inhaltsverzeichnis
	\begin{spacing}{1.1}
		% 2 = Kapitel bis Unterkapitel (1.1.1) im Inhaltsverzeichnis
		\setcounter{tocdepth}{2}
		\tableofcontents
	\end{spacing}
	\newpage

	\renewcommand{\thepage}{\arabic{page}}
	\setcounter{page}{1}

	% Inhalt -- neue Kapitel als Datei in content/ anlegen und hier eintragen
	%!TEX root = ../dokumentation.tex
% Schreibregel für alle Kapitel: EIN SATZ PRO ZEILE.
% Die PDF-Vorschau springt zeilengenau (SyncTeX); nur so landet sie genau auf
% dem Satz unter dem Cursor. Eine Leerzeile beginnt einen neuen Absatz.

\chapter{Einleitung}\label{ch:einleitung}

\section{Problemstellung}\label{sec:problemstellung}

Beschreibe hier, welches Problem die Arbeit löst und warum es relevant ist.
Jeder Satz steht in einer eigenen Zeile.

\section{Zielsetzung}\label{sec:zielsetzung}

Formuliere hier die Forschungsfrage oder das Ziel der Arbeit.

\section{Aufbau der Arbeit}\label{sec:aufbau}

\autoref{ch:grundlagen} erklärt die Grundlagen.
\autoref{ch:methodik} beschreibt das Vorgehen und \autoref{ch:umsetzung} die Umsetzung.
\autoref{ch:ergebnisse} stellt die Ergebnisse vor, \autoref{ch:fazit} fasst sie zusammen.

	%!TEX root = ../dokumentation.tex

\chapter{Grundlagen}\label{ch:grundlagen}

Dieses Kapitel erklärt die Begriffe und Verfahren, auf denen die Arbeit aufbaut.

\section{Erster Grundlagenbereich}\label{sec:grundlagen-erster}

Ein Satz, der eine Quelle braucht, bekommt sie direkt dahinter \citep{lamport1994latex}.

\section{Zweiter Grundlagenbereich}\label{sec:grundlagen-zweiter}

Platzhalter.

	%!TEX root = ../dokumentation.tex

\chapter{Methodik}\label{ch:methodik}

Dieses Kapitel beschreibt das Vorgehen und begründet die gewählten Methoden.

\section{Vorgehen}\label{sec:vorgehen}

Platzhalter.

\section{Bewertungskriterien}\label{sec:kriterien}

Platzhalter.

	%!TEX root = ../dokumentation.tex

\chapter{Umsetzung}\label{ch:umsetzung}

Dieses Kapitel beschreibt, wie das Konzept umgesetzt wurde.

\section{Architektur}\label{sec:architektur}

Platzhalter.

\section{Implementierung}\label{sec:implementierung}

Platzhalter.

	%!TEX root = ../dokumentation.tex

\chapter{Ergebnisse}\label{ch:ergebnisse}

Dieses Kapitel stellt die Ergebnisse vor und ordnet sie ein.

\section{Auswertung}\label{sec:auswertung}

Platzhalter.

\section{Diskussion}\label{sec:diskussion}

Platzhalter.

	%!TEX root = ../dokumentation.tex

\chapter{Fazit und Ausblick}\label{ch:fazit}

\section{Fazit}\label{sec:fazit}

Beantworte hier die Forschungsfrage aus \autoref{sec:zielsetzung}.

\section{Ausblick}\label{sec:ausblick}

Platzhalter.

	%!TEX root = ../dokumentation.tex
% SPICKZETTEL: jedes Element der Vorlage einmal. Im PDF steht das Ergebnis,
% hier der Quelltext dazu. Zum Entfernen die Zeile %!TEX root = ../dokumentation.tex
% SPICKZETTEL: jedes Element der Vorlage einmal. Im PDF steht das Ergebnis,
% hier der Quelltext dazu. Zum Entfernen die Zeile %!TEX root = ../dokumentation.tex
% SPICKZETTEL: jedes Element der Vorlage einmal. Im PDF steht das Ergebnis,
% hier der Quelltext dazu. Zum Entfernen die Zeile \input{content/beispiele}
% in dokumentation.tex löschen.

\chapter{Beispiele (Spickzettel)}\label{ch:beispiele}

Dieses Kapitel zeigt jedes Element der Vorlage einmal.
Den Quelltext dazu enthält \texttt{content/beispiele.tex}.

\section{Zitieren}\label{sec:bsp-zitieren}

Ein Beleg in Klammern steht am Satzende \citep{shannon1948mathematical}.
Mit Seitenangabe sieht er so aus \citep[S.~42]{knuth1984texbook}.
Im Satz zitiert: \citet{vaswani2017attention} stellen die Transformer-Architektur vor.
Online-Quellen tragen ihre URL im Literaturverzeichnis \citep{overleaf2026bibtex}.
Neue Quellen kommen in \texttt{literatur.bib}; das Literaturverzeichnis zeigt nur zitierte Einträge.

\section{Abkürzungen und Glossar}\label{sec:bsp-abkuerzungen}

Die erste Verwendung von \ac{API} schreibt die Langform aus, danach steht nur noch \ac{API}.
Glossarbegriffe wie \gls{synctex} werden mit dem Glossar verlinkt.

\section{Abbildungen}\label{sec:bsp-abbildungen}

\autoref{fig:bsp-bild} zeigt ein Bild; eigene Bilder liegen in \texttt{images/}.

\begin{figure}[htbp]
  \centering
  \includegraphics[width=0.5\textwidth]{example-image}
  \caption{Beispielbild: Dateiname ohne Endung angeben}
  \label{fig:bsp-bild}
\end{figure}

\section{Diagramme}\label{sec:bsp-diagramme}

\autoref{fig:bsp-diagramm} nutzt den TikZ-Stilbaukasten aus \texttt{ads/header.tex}.
Die Farbe aller Diagramme steuert \texttt{akzent} an einer Stelle.

\begin{figure}[htbp]
  \centering
  \begin{adjustbox}{max width=\textwidth}
  \begin{tikzpicture}[flow, node distance=6mm and 8mm]
    \node[fio]                   (in)   {Eingabe};
    \node[fstage, right=of in]   (prep) {Vorverarbeitung};
    \node[fgate,  right=of prep] (gate) {sicher?};
    \node[fkey,   right=of gate] (key)  {teurer\\Schritt};
    \node[fio,    right=of key]  (out)  {Ausgabe};
    \node[fmuted, below=of gate] (fast) {schneller Pfad};
    \draw[fedge] (in) -- (prep);
    \draw[fedge] (prep) -- (gate);
    \draw[fedge] (gate) -- node[flab, above] {nein} (key);
    \draw[fcond] (gate) -- node[flab, right] {ja} (fast);
    \draw[fedge] (key) -- (out);
    \draw[fcond] (fast) -| (out);
    \node[fnote, below=1mm of in] {Schritt 1};
  \end{tikzpicture}
  \end{adjustbox}
  \caption[Beispieldiagramm]{Beispieldiagramm mit den Stilen \texttt{fio}, \texttt{fstage}, \texttt{fgate}, \texttt{fkey} und \texttt{fmuted}. Der Kurztitel in eckigen Klammern erscheint im Abbildungsverzeichnis.}
  \label{fig:bsp-diagramm}
\end{figure}

\section{Tabellen}\label{sec:bsp-tabellen}

\autoref{tab:bsp-tabelle} nutzt \texttt{booktabs}: nur waagerechte Linien, die Beschriftung steht oben.

\begin{table}[htbp]
  \centering
  \caption{Beispieltabelle}
  \label{tab:bsp-tabelle}
  \begin{tabular}{lrr}
    \toprule
    Variante & Genauigkeit & Zeit pro Dokument \\
    \midrule
    Variante A & 72,4\,\% & 3\,ms \\
    Variante B & 88,0\,\% & 40\,ms \\
    \bottomrule
  \end{tabular}
\end{table}

\section{Quellcode}\label{sec:bsp-quellcode}

\autoref{lst:bsp-code} zeigt ein Listing; die Sprache stellt \texttt{[language=Java]} pro Listing um.

\begin{lstlisting}[caption={Beispiel-Listing in Python}, label={lst:bsp-code}]
def begruessung(name: str) -> str:
    # Umlaute im Code funktionieren: äöü ß
    return f"Hallo, {name}!"
\end{lstlisting}

\section{Formeln}\label{sec:bsp-formeln}

Nummerierte Formeln stehen in einer \texttt{equation}-Umgebung, \eqref{eq:bsp-f1} verweist darauf.
\begin{equation}
  F_1 = \frac{2 \cdot P \cdot R}{P + R}
  \label{eq:bsp-f1}
\end{equation}

\section{Sonstiges}\label{sec:bsp-sonstiges}

Eine Fußnote hängt direkt am Wort.\footnote{So sieht eine Fußnote aus.}
\enquote{Anführungszeichen} setzt \texttt{\textbackslash enquote} passend zur Sprache.
Querverweise mit \texttt{\textbackslash autoref} schreiben die Art des Ziels selbst davor: \autoref{sec:bsp-zitieren}, \autoref{fig:bsp-bild}, \autoref{tab:bsp-tabelle}.
Beträge mit Euro: 120\,\euro{}.

\begin{notiz}
Rote Entwurfsnotiz für dich selbst.
Vor der Abgabe blendet \texttt{\textbackslash excludecomment\{notiz\}} in \texttt{ads/header.tex} alle Notizen aus.
\end{notiz}

% in dokumentation.tex löschen.

\chapter{Beispiele (Spickzettel)}\label{ch:beispiele}

Dieses Kapitel zeigt jedes Element der Vorlage einmal.
Den Quelltext dazu enthält \texttt{content/beispiele.tex}.

\section{Zitieren}\label{sec:bsp-zitieren}

Ein Beleg in Klammern steht am Satzende \citep{shannon1948mathematical}.
Mit Seitenangabe sieht er so aus \citep[S.~42]{knuth1984texbook}.
Im Satz zitiert: \citet{vaswani2017attention} stellen die Transformer-Architektur vor.
Online-Quellen tragen ihre URL im Literaturverzeichnis \citep{overleaf2026bibtex}.
Neue Quellen kommen in \texttt{literatur.bib}; das Literaturverzeichnis zeigt nur zitierte Einträge.

\section{Abkürzungen und Glossar}\label{sec:bsp-abkuerzungen}

Die erste Verwendung von \ac{API} schreibt die Langform aus, danach steht nur noch \ac{API}.
Glossarbegriffe wie \gls{synctex} werden mit dem Glossar verlinkt.

\section{Abbildungen}\label{sec:bsp-abbildungen}

\autoref{fig:bsp-bild} zeigt ein Bild; eigene Bilder liegen in \texttt{images/}.

\begin{figure}[htbp]
  \centering
  \includegraphics[width=0.5\textwidth]{example-image}
  \caption{Beispielbild: Dateiname ohne Endung angeben}
  \label{fig:bsp-bild}
\end{figure}

\section{Diagramme}\label{sec:bsp-diagramme}

\autoref{fig:bsp-diagramm} nutzt den TikZ-Stilbaukasten aus \texttt{ads/header.tex}.
Die Farbe aller Diagramme steuert \texttt{akzent} an einer Stelle.

\begin{figure}[htbp]
  \centering
  \begin{adjustbox}{max width=\textwidth}
  \begin{tikzpicture}[flow, node distance=6mm and 8mm]
    \node[fio]                   (in)   {Eingabe};
    \node[fstage, right=of in]   (prep) {Vorverarbeitung};
    \node[fgate,  right=of prep] (gate) {sicher?};
    \node[fkey,   right=of gate] (key)  {teurer\\Schritt};
    \node[fio,    right=of key]  (out)  {Ausgabe};
    \node[fmuted, below=of gate] (fast) {schneller Pfad};
    \draw[fedge] (in) -- (prep);
    \draw[fedge] (prep) -- (gate);
    \draw[fedge] (gate) -- node[flab, above] {nein} (key);
    \draw[fcond] (gate) -- node[flab, right] {ja} (fast);
    \draw[fedge] (key) -- (out);
    \draw[fcond] (fast) -| (out);
    \node[fnote, below=1mm of in] {Schritt 1};
  \end{tikzpicture}
  \end{adjustbox}
  \caption[Beispieldiagramm]{Beispieldiagramm mit den Stilen \texttt{fio}, \texttt{fstage}, \texttt{fgate}, \texttt{fkey} und \texttt{fmuted}. Der Kurztitel in eckigen Klammern erscheint im Abbildungsverzeichnis.}
  \label{fig:bsp-diagramm}
\end{figure}

\section{Tabellen}\label{sec:bsp-tabellen}

\autoref{tab:bsp-tabelle} nutzt \texttt{booktabs}: nur waagerechte Linien, die Beschriftung steht oben.

\begin{table}[htbp]
  \centering
  \caption{Beispieltabelle}
  \label{tab:bsp-tabelle}
  \begin{tabular}{lrr}
    \toprule
    Variante & Genauigkeit & Zeit pro Dokument \\
    \midrule
    Variante A & 72,4\,\% & 3\,ms \\
    Variante B & 88,0\,\% & 40\,ms \\
    \bottomrule
  \end{tabular}
\end{table}

\section{Quellcode}\label{sec:bsp-quellcode}

\autoref{lst:bsp-code} zeigt ein Listing; die Sprache stellt \texttt{[language=Java]} pro Listing um.

\begin{lstlisting}[caption={Beispiel-Listing in Python}, label={lst:bsp-code}]
def begruessung(name: str) -> str:
    # Umlaute im Code funktionieren: äöü ß
    return f"Hallo, {name}!"
\end{lstlisting}

\section{Formeln}\label{sec:bsp-formeln}

Nummerierte Formeln stehen in einer \texttt{equation}-Umgebung, \eqref{eq:bsp-f1} verweist darauf.
\begin{equation}
  F_1 = \frac{2 \cdot P \cdot R}{P + R}
  \label{eq:bsp-f1}
\end{equation}

\section{Sonstiges}\label{sec:bsp-sonstiges}

Eine Fußnote hängt direkt am Wort.\footnote{So sieht eine Fußnote aus.}
\enquote{Anführungszeichen} setzt \texttt{\textbackslash enquote} passend zur Sprache.
Querverweise mit \texttt{\textbackslash autoref} schreiben die Art des Ziels selbst davor: \autoref{sec:bsp-zitieren}, \autoref{fig:bsp-bild}, \autoref{tab:bsp-tabelle}.
Beträge mit Euro: 120\,\euro{}.

\begin{notiz}
Rote Entwurfsnotiz für dich selbst.
Vor der Abgabe blendet \texttt{\textbackslash excludecomment\{notiz\}} in \texttt{ads/header.tex} alle Notizen aus.
\end{notiz}

% in dokumentation.tex löschen.

\chapter{Beispiele (Spickzettel)}\label{ch:beispiele}

Dieses Kapitel zeigt jedes Element der Vorlage einmal.
Den Quelltext dazu enthält \texttt{content/beispiele.tex}.

\section{Zitieren}\label{sec:bsp-zitieren}

Ein Beleg in Klammern steht am Satzende \citep{shannon1948mathematical}.
Mit Seitenangabe sieht er so aus \citep[S.~42]{knuth1984texbook}.
Im Satz zitiert: \citet{vaswani2017attention} stellen die Transformer-Architektur vor.
Online-Quellen tragen ihre URL im Literaturverzeichnis \citep{overleaf2026bibtex}.
Neue Quellen kommen in \texttt{literatur.bib}; das Literaturverzeichnis zeigt nur zitierte Einträge.

\section{Abkürzungen und Glossar}\label{sec:bsp-abkuerzungen}

Die erste Verwendung von \ac{API} schreibt die Langform aus, danach steht nur noch \ac{API}.
Glossarbegriffe wie \gls{synctex} werden mit dem Glossar verlinkt.

\section{Abbildungen}\label{sec:bsp-abbildungen}

\autoref{fig:bsp-bild} zeigt ein Bild; eigene Bilder liegen in \texttt{images/}.

\begin{figure}[htbp]
  \centering
  \includegraphics[width=0.5\textwidth]{example-image}
  \caption{Beispielbild: Dateiname ohne Endung angeben}
  \label{fig:bsp-bild}
\end{figure}

\section{Diagramme}\label{sec:bsp-diagramme}

\autoref{fig:bsp-diagramm} nutzt den TikZ-Stilbaukasten aus \texttt{ads/header.tex}.
Die Farbe aller Diagramme steuert \texttt{akzent} an einer Stelle.

\begin{figure}[htbp]
  \centering
  \begin{adjustbox}{max width=\textwidth}
  \begin{tikzpicture}[flow, node distance=6mm and 8mm]
    \node[fio]                   (in)   {Eingabe};
    \node[fstage, right=of in]   (prep) {Vorverarbeitung};
    \node[fgate,  right=of prep] (gate) {sicher?};
    \node[fkey,   right=of gate] (key)  {teurer\\Schritt};
    \node[fio,    right=of key]  (out)  {Ausgabe};
    \node[fmuted, below=of gate] (fast) {schneller Pfad};
    \draw[fedge] (in) -- (prep);
    \draw[fedge] (prep) -- (gate);
    \draw[fedge] (gate) -- node[flab, above] {nein} (key);
    \draw[fcond] (gate) -- node[flab, right] {ja} (fast);
    \draw[fedge] (key) -- (out);
    \draw[fcond] (fast) -| (out);
    \node[fnote, below=1mm of in] {Schritt 1};
  \end{tikzpicture}
  \end{adjustbox}
  \caption[Beispieldiagramm]{Beispieldiagramm mit den Stilen \texttt{fio}, \texttt{fstage}, \texttt{fgate}, \texttt{fkey} und \texttt{fmuted}. Der Kurztitel in eckigen Klammern erscheint im Abbildungsverzeichnis.}
  \label{fig:bsp-diagramm}
\end{figure}

\section{Tabellen}\label{sec:bsp-tabellen}

\autoref{tab:bsp-tabelle} nutzt \texttt{booktabs}: nur waagerechte Linien, die Beschriftung steht oben.

\begin{table}[htbp]
  \centering
  \caption{Beispieltabelle}
  \label{tab:bsp-tabelle}
  \begin{tabular}{lrr}
    \toprule
    Variante & Genauigkeit & Zeit pro Dokument \\
    \midrule
    Variante A & 72,4\,\% & 3\,ms \\
    Variante B & 88,0\,\% & 40\,ms \\
    \bottomrule
  \end{tabular}
\end{table}

\section{Quellcode}\label{sec:bsp-quellcode}

\autoref{lst:bsp-code} zeigt ein Listing; die Sprache stellt \texttt{[language=Java]} pro Listing um.

\begin{lstlisting}[caption={Beispiel-Listing in Python}, label={lst:bsp-code}]
def begruessung(name: str) -> str:
    # Umlaute im Code funktionieren: äöü ß
    return f"Hallo, {name}!"
\end{lstlisting}

\section{Formeln}\label{sec:bsp-formeln}

Nummerierte Formeln stehen in einer \texttt{equation}-Umgebung, \eqref{eq:bsp-f1} verweist darauf.
\begin{equation}
  F_1 = \frac{2 \cdot P \cdot R}{P + R}
  \label{eq:bsp-f1}
\end{equation}

\section{Sonstiges}\label{sec:bsp-sonstiges}

Eine Fußnote hängt direkt am Wort.\footnote{So sieht eine Fußnote aus.}
\enquote{Anführungszeichen} setzt \texttt{\textbackslash enquote} passend zur Sprache.
Querverweise mit \texttt{\textbackslash autoref} schreiben die Art des Ziels selbst davor: \autoref{sec:bsp-zitieren}, \autoref{fig:bsp-bild}, \autoref{tab:bsp-tabelle}.
Beträge mit Euro: 120\,\euro{}.

\begin{notiz}
Rote Entwurfsnotiz für dich selbst.
Vor der Abgabe blendet \texttt{\textbackslash excludecomment\{notiz\}} in \texttt{ads/header.tex} alle Notizen aus.
\end{notiz}
 % Spickzettel mit allen Elementen -- Zeile löschen, wenn nicht mehr gebraucht

	% Anhang: \appendix steht in der Datei, die Kapitel heißen dann A, B, ...
	%!TEX root = ../dokumentation.tex
% Ab \appendix werden Kapitel mit Buchstaben nummeriert (Anhang A, B, ...).

\appendix

\chapter{Zusatzmaterial}\label{app:zusatz}

Hierher gehört, was den Lesefluss stören würde: lange Tabellen, vollständiger Quellcode, Fragebögen.


	% Verzeichnisse
	\clearpage
	\pagenumbering{roman}

	% Abbildungsverzeichnis
	\cleardoublepage
	\phantomsection \label{listoffig}
	\addcontentsline{toc}{chapter}{\listfigurename}
	\listoffigures

	% Tabellenverzeichnis
	\cleardoublepage
	\phantomsection \label{listoftab}
	\addcontentsline{toc}{chapter}{\listtablename}
	\listoftables

	% Quellcodeverzeichnis (optional)
	%\cleardoublepage
	%\phantomsection \label{listoflst}
	%\addcontentsline{toc}{chapter}{\lstlistlistingname}
	%\lstlistoflistings

	% Literaturverzeichnis
	\cleardoublepage
	\phantomsection \label{listoflit}
	\addcontentsline{toc}{chapter}{\bibname}

	% Stil: Harvard (Autor, Jahr); agsm-de.bst ist agsm mit deutschen Wörtern
	% (S., Hrsg., Aufl.). Andere Stile, die mit harvard.sty funktionieren
	% (alle englisch): dcu, kluwer, jmr, apsr, nederlands, jphysicsB.
	\ifenglisch\bibliographystyle{agsm}\else\bibliographystyle{agsm-de}\fi
	\bibliography{literatur}

	% Abkürzungsverzeichnis (zeigt nur Abkürzungen, die im Text per \ac{} vorkommen)
	\cleardoublepage
	\phantomsection \label{listofacs}
	\addcontentsline{toc}{chapter}{\abkuerzungstitel}
	%!TEX root = ../dokumentation.tex

\chapter*{\abkuerzungstitel}
% Im Text: \ac{API} schreibt beim ersten Mal die Langform (als Fußnote), danach
% nur die Abkürzung. Gedruckt werden nur Abkürzungen, die im Text per \ac{}
% vorkommen (printonlyused); ohne jedes \ac{} bleibt die Liste leer.
% [WYSIWYG] = längste Abkürzung, bestimmt die Spaltenbreite.
\begin{acronym}[WYSIWYG]
\setlength{\itemsep}{-\parsep}

\acro{API}{Application Programming Interface}
\acro{DHBW}{Duale Hochschule Baden-Württemberg}
\acro{PDF}{Portable Document Format}
\end{acronym}


	% Glossar
	% \glsaddall druckt alle Einträge aus ads/glossary.tex, auch die im Text
	% nicht per \gls{} verlinkten.
	\glsaddall
	\printglossary[style=altlist,title=\glossartitel]
\end{document}
